% SPDX-License-Identifier: CC-BY-4.0
% Copyright 2025 Florian Aram Feuerriegel — kassensturz.org
% Abschnitt 5 — Validierung (tests/)

\section{Validierung}
\label{sec:validierung}

Ein vereinfachtes Modell kann ökonomisch unsinnig rechnen, ohne dass es
auffällt: Jede Einzelformel ist plausibel, das Zusammenspiel nicht. Zwei
Fachprüfungen im September 2026 fanden in früheren Fassungen Fehler dieser
Art, die keiner der damaligen Tests erfassen konnte. Zinsen wurden in der
Schuldenfortschreibung ein zweites Mal gezählt, der Primärsaldo war um den
Faktor 100 verschoben, und eine Senkung des Rentenbeitrags verbesserte den
Saldo und machte alle Haushalte reicher, weil die gekürzten Renten bei keinem
Haushalt fehlten. Ein Test, der nur den Ist-Zustand festschreibt, bemerkt so
etwas nicht. Die Tests in \datei{tests/} prüfen deshalb Eigenschaften, die
jede zulässige Politik erfüllen muss:

\begin{description}
\item[Normierung.] Im Status quo ist das Arbeitsangebot unverändert,
  $\lambda_i = 1$, und Inzidenz und Nachfragelücke sind null. Kennzahlen des
  Staates und Nachfrage sind damit Differenzen zum Status quo derselben Periode.
\item[Buchungsidentität.] Die Summe über die Haushalte entspricht der direkten
  Buchung des Staates, ohne die Rückwirkung über die Nachfrage: exakt für
  Kinder- und Bürgergeld, Renten, \COz-Last und Klimageld, Unternehmen- und
  Vermögensteuern; auf 15\,\% genau für Umsatzsteuer, Beitragssätze und
  Beitragsbemessungsgrenze, bei denen Staat und Haushalte verschiedene
  Bemessungsgrundlagen haben.
\item[Keine dominante Strategie.] Keine Stellgröße, geprüft an elf Werten über
  ihren Spielraum, und kein Paar von Stellgrößen am Tisch an den Rändern ihres
  Spielraums verbessert in der letzten Periode den Saldo und das Einkommen aller
  Haushaltstypen, ohne Gini, Emissionen oder BIP-Niveau zu verschlechtern.
\item[Komparative Statik.] Die Vorzeichen der Reaktionen entsprechen der
  Theorie: Ein höherer \COz-Preis senkt die Emissionen, eine höhere
  Körperschaftsteuer die Investitionen, ein höherer Eingangssatz das
  Arbeitsangebot, ein Spitzensatz über 45\,\% das Einkommen im obersten Prozent.
  Konsolidierung und Entlastung wirken über die Nachfrage betragsgleich mit
  umgekehrtem Vorzeichen, und Steuerpolitik verschiebt das Niveau des BIP, nicht
  seine Wachstumsrate.
\item[Dimensionen.] Zinslast in Prozentpunkten des BIP, ein Zinssatz für
  Haushalt und Schuldendynamik, keine doppelte Zählung von Zinsen, konstante
  Einnahmenquote im Status quo.
\item[Invarianten.] Alle Kennzahlen sind endlich und liegen in ihrem
  Definitionsbereich, für alle Voreinstellungen, Extremwerte der Stellgrößen
  und Pfade bis zu zwölf Perioden.
\item[Kalibrierung.] Aufkommen je Steuerart, Gini und Palma liegen in der
  Größenordnung der amtlichen Werte (Tabelle~\ref{tab:kalibrierung}).
\item[Tarif.] Der Tarif weicht an keiner Stelle um mehr als 1\,\% von §\,32a
  EStG 2026 ab, ist stetig und monoton, und der Grenzsteuersatz ist seine
  Ableitung.
\item[Dokumentation.] Jede Zahl in Text, Tabellen und Abbildungen entspricht
  dem Modell, und jede Stellgröße am Tisch ist gemessen.
\end{description}

Die Buchungsidentität ist die wichtigste dieser Eigenschaften. Nur wenn der
Staat bucht, was bei den Haushalten ankommt, sind Haushalts- und
Verteilungswirkung derselben Politik vergleichbar. Die zweite ist das Fehlen
einer dominanten Strategie: Ein Planspiel über Zielkonflikte, in dem eine
Stellgröße alle Ziele zugleich verbessert, lehrt das Falsche. Findet der Test
eine solche Politik, fehlt dem Modell ein Kanal, über den sie etwas kostet.
